% Paper 4, §8 — Are SAE features canonical?
% Thesis: the decomposition uniqueness question. Leask 2025 (ICLR) gives
% two pieces of counter-evidence to the Bricken/Templeton "true features"
% framing — meta-SAEs (non-atomicity) and SAE stitching (incompleteness).
% Transcoders (Dunefsky 2024) are the structural cousin: they reproduce
% the *function* of an MLP sublayer rather than its activations, and
% Lindsey 2025's cross-layer transcoder is the architectural primitive
% of §11's circuit tracing — but they inherit the canonicality concern.

\section{Are SAE features canonical?}
\label{sec:sae-canonicality}
\label{sec:sae-canonical}

The architectures of \S\ref{sec:sparse-autoencoders} give us a working
pipeline for converting polysemantic neuron-basis activations into a
sparse, named \glsterm{feature}{feature} basis at scale. They do not, on their own, tell
us whether that basis is \emph{the} basis. The implicit framing of the
production-scale \glsterm{sparse-autoencoder}{SAE} program --- that scaling $M$ and pushing further
down the $L_0$-vs-reconstruction Pareto curve converges toward a
canonical decomposition of model activations into atomic, interpretable
features --- is a load-bearing assumption: if features are
\emph{canonical}, then \glsterm{sparse-autoencoder}{SAE} features are the units of the model's
internal computation, and downstream interpretability claims (probes
cite them, patching targets them, circuits compose them) inherit
their canonicality. \citet{leask2025-sae-not-canonical} make a direct
empirical argument that this assumption is unwarranted, with two
techniques that \glsterm{probe-2}{probe} atomicity and completeness separately. We work
through the argument, then take up \emph{transcoders}
\citep{dunefsky2024-transcoders} --- structural cousins of SAEs that
sit at the input-output boundary of an MLP rather than on the
activation itself --- and show that the additional analytical leverage
they provide for \glsterm{circuit}{circuit} work in \S\ref{sec:circuits} comes with the
same canonicality concern inherited unchanged.

\subsection{What ``canonical'' would mean}
\label{sec:canonical-definition}

\citet[\S 1]{leask2025-sae-not-canonical} fix the terminology before
attacking the claim\footnote{KB excerpt:
\texttt{kb/excerpts/leask2025-\glsterm{sparse-autoencoder}{sae}-not-canonical.md\#sec-1-canonical}.}.
A set of units of analysis is \textbf{canonical} when it is
simultaneously: \textbf{unique} --- the same set is recovered (modulo
a known equivalence such as rotation or permutation) by independent
training runs and across reasonable choices of dictionary size $M$;
\textbf{complete} --- every \glsterm{feature}{feature} the model uses is represented;
and \textbf{atomic} --- each unit is irreducible, not decomposable into
smaller units in the same dictionary family. The
\citet{bricken2023-monosemanticity} ``Towards \glsterm{monosemanticity}{Monosemanticity}'' paper
and the \citet{templeton2024-scaling-monosemanticity} scaling
follow-up both lean on a strong-form version of this picture: \glsterm{sparse-autoencoder}{SAE}
features are the \emph{true} units, and a sufficiently scaled \glsterm{sparse-autoencoder}{SAE}
recovers them. Both citations are HTML-only on
\texttt{\glsterm{transformer}{transformer}-circuits.pub}; we cite them by paper-\glsterm{key}{key} with
attribution pending Phase~2 transcription
\deepencite{bricken2023-monosemanticity §2, citation-pending}
\deepencite{templeton2024-scaling-monosemanticity §1, citation-pending}.

The empirical content of the canonicality claim is testable. If
features are canonical, then (i) a wider \glsterm{sparse-autoencoder}{SAE} on the same activations
should not contain features absent from a sufficiently wide narrower
\glsterm{sparse-autoencoder}{SAE} --- it should refine, not enlarge, the inventory; and (ii) a
\glsterm{feature}{feature} in a wide \glsterm{sparse-autoencoder}{SAE} should not decompose into a sparse combination
of features from a narrower \glsterm{sparse-autoencoder}{SAE} --- it should be a single unit. The
two techniques of \citet{leask2025-sae-not-canonical} test exactly
these two predictions, and report failure on both.

\subsection{Meta-SAEs: large-SAE features are not atomic}
\label{sec:meta-saes}

The first technique trains an \glsterm{sparse-autoencoder}{SAE} on the \emph{decoder matrix} of a
target \glsterm{sparse-autoencoder}{SAE} rather than on model activations. Each row of the
\glsterm{meta-sae}{meta-SAE}'s input is a single decoder column $\mathbf{d}_i \in
\mathbb{R}^n$ from the target \glsterm{sparse-autoencoder}{SAE}; the \glsterm{meta-sae}{meta-SAE} learns a sparse
dictionary of \emph{meta-latents} into which the target's \glsterm{feature}{feature}
directions decompose. The reported finding\footnote{KB excerpt:
\texttt{kb/excerpts/leask2025-\glsterm{sparse-autoencoder}{sae}-not-canonical.md\#sec-1-einstein}.}
is the now-canonical example:

\begin{quote}
\itshape
[Figure 1] Decomposition of an \glsterm{sparse-autoencoder}{SAE} latent representing ``Einstein''
into a set of interpretable meta-latents. \dots\ revealing how a
complex concept can be represented as a sparse combination of
meta-latents.~\citep[\S 1]{leask2025-sae-not-canonical}
\end{quote}

\noindent The Einstein latent in the target \glsterm{sparse-autoencoder}{SAE} decomposes into
meta-latents corresponding to ``scientist'', ``Germany'', ``famous
person'', and a ``starts with E-'' \glsterm{feature}{feature}. Atomicity --- the claim
that the Einstein direction is irreducible within the \glsterm{sparse-autoencoder}{SAE}'s \glsterm{feature}{feature}
family --- fails on inspection: it is a sparse combination of more
elementary directions discovered without looking at any new model
activations, just at the target \glsterm{sparse-autoencoder}{SAE}'s own decoder. The summary
characterization is direct\footnote{KB excerpt:
\texttt{kb/excerpts/leask2025-\glsterm{sparse-autoencoder}{sae}-not-canonical.md\#abstract}.}:

\begin{quote}
\itshape
Using meta-SAEs --- SAEs trained on the decoder matrix of another
\glsterm{sparse-autoencoder}{SAE} --- we find that latents in SAEs often decompose into
combinations of latents from a smaller \glsterm{sparse-autoencoder}{SAE}, showing that larger \glsterm{sparse-autoencoder}{SAE}
latents are not atomic.~\citep[abstract]{leask2025-sae-not-canonical}
\end{quote}

\noindent The decoder-matrix substrate is the load-bearing methodological
choice: meta-SAEs do not introduce new activation data, so the
decomposition of $\mathbf{d}_\text{Einstein}$ into meta-latent
directions is a property of the target \glsterm{sparse-autoencoder}{SAE}'s learned geometry alone.
Whatever ``Einstein'' meant inside the target \glsterm{sparse-autoencoder}{SAE}, it meant it as a
sparse combination of more-elementary directions in
$\mathbb{R}^n$ --- and the \glsterm{meta-sae}{meta-SAE} recovers them.

\subsection{SAE stitching: smaller SAEs are incomplete}
\label{sec:sae-stitching}

The second technique tests completeness by inserting or swapping
latents from a larger \glsterm{sparse-autoencoder}{SAE} into a smaller one and measuring
reconstruction performance on held-out activations. The paper's
classification\footnote{KB excerpt:
\texttt{kb/excerpts/leask2025-\glsterm{sparse-autoencoder}{sae}-not-canonical.md\#abstract}.} sorts
larger-\glsterm{sparse-autoencoder}{SAE} latents into two disjoint sets:

\begin{itemize}
  \item \textbf{Reconstruction latents.} When inserted into the
    smaller \glsterm{sparse-autoencoder}{SAE}, they can replace --- with comparable behavior ---
    one or more latents already present. They are coordinate
    refinements of features the smaller \glsterm{sparse-autoencoder}{SAE} already had.
  \item \textbf{Novel latents.} When inserted, they
    \emph{improve} the smaller \glsterm{sparse-autoencoder}{SAE}'s reconstruction. They are
    features the smaller \glsterm{sparse-autoencoder}{SAE} did not have; absent in its dictionary;
    inaccessible to its forward pass.
\end{itemize}

\noindent The empirical finding is that novel latents exist in
non-trivial numbers: scaling $M$ does not just refine, it discovers.
The methodological consequence
\citep[\S 1]{leask2025-sae-not-canonical} is that any single-width \glsterm{sparse-autoencoder}{SAE}
is provably incomplete relative to a sufficiently wider one --- and
since the field has no a-priori stopping criterion for $M$, the
``complete dictionary'' premise of the canonicality claim has no
operational test that any finite \glsterm{sparse-autoencoder}{SAE} can pass. The contrast with
\citet{bricken2023-monosemanticity}'s earlier \emph{\glsterm{feature}{feature}-splitting}
framing is sharp:
\citeauthor{leask2025-sae-not-canonical} note that under \glsterm{feature}{feature}
splitting alone, ``latents from smaller SAEs `split' into multiple,
more fine-grained latents in larger
SAEs''~\citep[\S 1]{leask2025-sae-not-canonical}\footnote{KB excerpt:
\texttt{kb/excerpts/leask2025-\glsterm{sparse-autoencoder}{sae}-not-canonical.md\#sec-1-\glsterm{feature}{feature}-splitting}.}
--- a relationship in which the smaller dictionary is \emph{coarser}
but still complete. Stitching's novel latents fail this picture: some
larger-\glsterm{sparse-autoencoder}{SAE} features have no coarse-grained counterpart in the smaller
dictionary at all.

\begin{contradiction}
\textbf{Are \glsterm{sparse-autoencoder}{SAE} features canonical?} Two lines of evidence strain the
strong-form ``true features'' framing implicit in
\citet{bricken2023-monosemanticity} and
\citet{templeton2024-scaling-monosemanticity}.
\textsf{(1)~Non-atomicity.} Meta-SAEs trained on the decoder matrix
$\mathbf{W}_\text{dec}$ of a target \glsterm{sparse-autoencoder}{SAE} recover a sparse decomposition
of the target's \glsterm{feature}{feature} directions; an ``Einstein'' latent decomposes
into ``scientist'', ``Germany'', ``famous person''
\citep[\S 1]{leask2025-sae-not-canonical}. The target's atoms are not
atoms.
\textsf{(2)~Incompleteness.} \glsterm{sae-stitching}{SAE stitching} shows that larger SAEs
contain \emph{novel latents} --- features absent from smaller SAEs
that improve reconstruction when transplanted in
\citep[abstract]{leask2025-sae-not-canonical}. No fixed-width \glsterm{sparse-autoencoder}{SAE}
qualifies as complete relative to a sufficiently wider trained
counterpart.
\textbf{What evidence would close it.} A basis-finding criterion under
which two independent \glsterm{sparse-autoencoder}{SAE} training runs --- at different widths
$M_1, M_2$, different sparsity targets, different random seeds, and
on the same activations --- converge on the same \glsterm{feature}{feature} inventory
modulo a known equivalence (rotation, sign, permutation, scale).
Until that criterion is exhibited and shown to hold empirically, the
field's working assumption is the pragmatic one
\citet[\S 1]{leask2025-sae-not-canonical} themselves
recommend\footnote{KB excerpt:
\texttt{kb/excerpts/leask2025-\glsterm{sparse-autoencoder}{sae}-not-canonical.md\#sec-1-still-useful}.}:
``the choice of dictionary size is subjective. We suggest taking a
pragmatic approach to applying SAEs to \glsterm{mechanistic-interpretability}{mechanistic interpretability}
tasks, trying SAEs of several widths to see which is best
suited.''~\citep[\S 1]{leask2025-sae-not-canonical}
\end{contradiction}

\subsection{Transcoders: function approximation, not activation
decomposition}
\label{sec:transcoder-canonicality}

The architectural cousin of an \glsterm{sparse-autoencoder}{SAE} relevant to \glsterm{circuit}{circuit} analysis ---
introduced in \S\ref{sec:sae-transcoders} and reused as the primitive
of \S\ref{sec:circuits} --- is the \emph{\glsterm{transcoder}{transcoder}}
\citep{dunefsky2024-transcoders}. We restate the loss for reference,
because the canonicality picture rides on its faithfulness term being
defined against the MLP \emph{function} rather than against its
activation\footnote{KB excerpt:
\texttt{kb/excerpts/dunefsky2024-transcoders.md\#sec-3-1}.}:
\begin{equation}
  \mathcal{L}_{\text{TC}}(\mathbf{x})
  \;=\;
  \underbrace{\|\mathrm{MLP}(\mathbf{x}) - \mathrm{TC}(\mathbf{x})\|_2^2}_{\text{faithfulness}}
  \;+\;
  \lambda_1\,\|\mathbf{z}_{\text{TC}}(\mathbf{x})\|_1
  \label{eq:transcoder-loss-canonical}
\end{equation}
with $\mathrm{TC}(\mathbf{x}) = \mathbf{W}_\text{dec}\,\mathrm{ReLU}(\mathbf{W}_\text{enc}\mathbf{x} + \mathbf{b}_\text{enc}) + \mathbf{b}_\text{dec}$
from \citet[\S 3.1, Eq.~3--5]{dunefsky2024-transcoders}. The
contrast with an \glsterm{sparse-autoencoder}{SAE} on the MLP output is one substitution: an \glsterm{sparse-autoencoder}{SAE}
minimizes $\|\mathrm{MLP}(\mathbf{x}) - \hat{\mathbf{x}}\|_2^2$ where
$\hat{\mathbf{x}}$ depends on the \emph{post}-MLP activation; a
\glsterm{transcoder}{transcoder}'s $\mathrm{TC}(\mathbf{x})$ depends only on the
\emph{pre}-MLP input $\mathbf{x}$. The \glsterm{transcoder}{transcoder}'s \glsterm{feature}{feature} dictionary
is a decomposition of the MLP \emph{function}, not of any particular
output \citep[\S 1]{dunefsky2024-transcoders}\footnote{KB excerpt:
\texttt{kb/excerpts/dunefsky2024-transcoders.md\#sec-1-comparison}.}.

The methodological payoff
\citep[\S 1]{dunefsky2024-transcoders}\footnote{KB excerpt:
\texttt{kb/excerpts/dunefsky2024-transcoders.md\#sec-1-\glsterm{input-invariance}{input-invariance}}.}
is the \textbf{\glsterm{input-invariance}{input-invariance} factorization}. The connection from a
pre-MLP \glsterm{feature}{feature} $f_i$ (associated with decoder column $\mathbf{d}_i$)
to a post-MLP \glsterm{transcoder}{transcoder} \glsterm{feature}{feature} $f_j$ (associated with encoder row
$\mathbf{e}_j$) factors into an input-dependent term --- the
activations $f_i(\mathbf{x})$ and $f_j(\mathbf{x})$ at the input of
interest --- and an input-invariant term, the bilinear product
\begin{equation}
  W_{j i} \;:=\; \mathbf{e}_j^\top \mathbf{d}_i,
  \label{eq:transcoder-bilinear}
\end{equation}
which holds for any input. The motivating example
\citet[\S 1]{dunefsky2024-transcoders} give for why this matters is
worth quoting in full:

\begin{quote}
\itshape
SAEs cannot tell us about the general input-output behavior of an MLP
across all inputs. \dots\ Doing e.g. patching on one input shows that a
pre-MLP \glsterm{feature}{feature} for Polish last names is important for causing the
post-MLP \glsterm{feature}{feature} to activate. But on other inputs, would features other
than the Polish last name \glsterm{feature}{feature} also cause the post-MLP \glsterm{feature}{feature} to
fire (e.g. an English last names \glsterm{feature}{feature})? Could there be other inputs
where the Polish last names \glsterm{feature}{feature} fires but the post-MLP \glsterm{feature}{feature} does
not? We can see that without \glsterm{input-invariance}{input-invariance}, it is difficult to make
general claims about model
behavior.~\citep[\S 1]{dunefsky2024-transcoders}
\end{quote}

\noindent An \glsterm{sparse-autoencoder}{SAE}-on-output answers ``which features explain
\emph{this} activation?''; a \glsterm{transcoder}{transcoder} answers ``which features does
this MLP \emph{compute}, and how do they depend on its inputs?''. The
latter is the question \glsterm{circuit}{circuit} analysis needs answered, and the
\glsterm{cross-layer-transcoder}{cross-layer transcoder} of \citet{lindsey2025-circuit-tracing}
\deepencite{lindsey2025-circuit-tracing §3, citation-pending}
generalizes the bilinear factorization \eqref{eq:transcoder-bilinear}
to skip-layer \glsterm{feature}{feature} flow. That generalization is the architectural
primitive on which the attribution-graph construction of
\S\ref{sec:circuit-attribution} is built.

\begin{intuition}
A \glsterm{transcoder}{transcoder} is to an MLP what a sparse Taylor expansion is to a
smooth function: the original is a dense, hard-to-read black box; the
\glsterm{transcoder}{transcoder} is a sparse linear combination of basis functions
$\{\mathbf{d}_i\}$ that approximates it everywhere. The
\glsterm{input-invariance}{input-invariance} factorization is exactly the statement that
``coefficients of the expansion'' (the bilinear products
$\mathbf{e}_j^\top \mathbf{d}_i$) are properties of the function, while
``where the expansion is being evaluated'' (the activations
$f_i(\mathbf{x}), f_j(\mathbf{x})$) are properties of the input.
Returning to canonical math: the Jacobian of a post-\glsterm{transcoder}{transcoder} \glsterm{feature}{feature}
$f_j$ with respect to a pre-\glsterm{transcoder}{transcoder} \glsterm{feature}{feature} $f_i$ separates into a
bilinear weight term and a pointwise activation derivative ---
the structural condition that makes attribution graphs computable in
closed form rather than by sampling.
\end{intuition}

\subsection{Inheritance: transcoders carry the SAE canonicality
problem}
\label{sec:transcoder-inheritance}

The architectural difference between Eq.~\eqref{eq:transcoder-loss-canonical}
and the \glsterm{sparse-autoencoder}{SAE} loss of Eq.~\eqref{eq:sae-loss} does not, on its own,
resolve the canonicality concerns of
\S\S\ref{sec:meta-saes}--\ref{sec:sae-stitching}. A \glsterm{transcoder}{transcoder} is still
a sparse-dictionary architecture; its \glsterm{feature}{feature} directions
$\{\mathbf{d}_i\}$ are still trained against a reconstruction
objective with a sparsity penalty; nothing in
Eq.~\eqref{eq:transcoder-loss-canonical} privileges any particular
basis modulo the rotational/scale freedoms of dictionary learning.
Two independent transcoders trained on the same MLP at different
widths $d_\text{features}$ should be expected --- by the same
arguments \citet{leask2025-sae-not-canonical} make for SAEs --- to
exhibit incompleteness (the wider \glsterm{transcoder}{transcoder} finds function-features
the narrower one missed) and non-atomicity (a wider \glsterm{feature}{feature}
$\mathbf{d}_i^\text{wide}$ admits a sparse decomposition into
function-features of a narrower \glsterm{transcoder}{transcoder}). The
\citet{leask2025-sae-not-canonical} stitching protocol applied between
transcoders rather than between SAEs has not, as of 2026, been
published as a separate experiment, and we flag the inheritance as
empirically unconfirmed but theoretically expected
\deepencite{leask2025-sae-not-canonical §1, transcoder-extension-pending}.

The pragmatic consequence for \glsterm{circuit}{circuit} work is real. The bilinear
factorization of Eq.~\eqref{eq:transcoder-bilinear} is well-defined
\emph{within} a fixed \glsterm{transcoder}{transcoder}: given the
$(\mathbf{W}_\text{enc}, \mathbf{W}_\text{dec})$ trained for a
specific MLP, the input-invariant attribution weights $W_{ji}$ are
deterministic. But the \emph{interpretation} of the post-\glsterm{transcoder}{transcoder}
\glsterm{feature}{feature} $f_j$ --- what it ``represents'' --- is identified at the
basis level, and the basis is not canonical. Two attribution graphs
computed against two different transcoders trained on the same MLP
will produce different node sets, different edge weights, and
potentially different qualitative narratives about the underlying
\glsterm{circuit}{circuit}. This is the \emph{\glsterm{transcoder}{transcoder}-canonicality} concern inherited
from the \glsterm{sparse-autoencoder}{SAE} setting, and the cross-method synthesis of
\S\ref{sec:cross-method} returns to it: when \glsterm{activation-patching}{activation patching} at
the head/MLP level and attribution-graph tracing at the \glsterm{transcoder}{transcoder}-\glsterm{feature}{feature}
level disagree, one possibility is that the patching result is wrong;
another is that the \glsterm{attribution-graph}{attribution graph} reads features as canonical that
the model does not treat as canonical.

\subsection{Where the canonicality question stands}
\label{sec:canonicality-frontier}

The 2026 community position, as registered by
\citet[\S 1]{leask2025-sae-not-canonical} and consistent with
Anthropic's own move from ``true features''
\deepencite{templeton2024-scaling-monosemanticity §1, citation-pending}
to \glsterm{circuit}{circuit}-tracing's pragmatic \glsterm{transcoder}{transcoder} usage
\deepencite{lindsey2025-circuit-tracing §1, citation-pending}, is
\emph{pragmatic}: SAEs and transcoders are useful, basis-dependent
tools for converting polysemantic activations into legible sparse
representations, and the right working response to the canonicality
critique is to vary $M$, retrain, and check that load-bearing
interpretability claims survive across widths and seeds. This is the
recommendation \citeauthor{leask2025-sae-not-canonical} themselves
make: ``trying SAEs of several widths to see which is best
suited''~\citep[\S 1]{leask2025-sae-not-canonical}.

The \emph{theoretical} question --- whether canonical units exist at
all, and whether some other dictionary-learning architecture (one
explicitly designed to converge across widths) could find them ---
is open. The shape of an answer would be: a basis-finding criterion
that two independent runs satisfy modulo a known equivalence, and an
architecture for which the criterion is provably (or at least
robustly empirically) achievable. Neither is in hand as of 2026.

\subsection{Forward link}
\label{sec:canonicality-forward}

The \glsterm{transcoder}{transcoder} framing of \S\ref{sec:transcoder-canonicality} is the
architectural primitive of \S\ref{sec:circuits}, where
\citet{lindsey2025-circuit-tracing}'s cross-layer transcoders provide
the bilinear-attribution substrate
\eqref{eq:transcoder-bilinear} on which Anthropic's 2025 attribution
graphs on \glsterm{claude}{Claude} 3.5 Haiku are built
\deepencite{lindsey2025-circuit-tracing §3, citation-pending}. The
canonicality concern propagates forward: any \glsterm{circuit}{circuit} claim made
against a \glsterm{transcoder}{transcoder}-derived \glsterm{attribution-graph}{attribution graph} is a claim about
\emph{this} \glsterm{transcoder}{transcoder}'s features, not about a model-level canonical
basis. \S\ref{sec:cross-method} takes up the cross-method synthesis
question --- when two interpretability methods agree on a phenomenon
(the IOI \glsterm{circuit}{circuit}, factual recall localization), and when they
disagree (the \glsterm{truth-direction-2}{truth direction} across datasets, \glsterm{sparse-autoencoder}{SAE} \glsterm{feature}{feature} inventories
across widths) --- and uses the canonicality picture from this
section to explain why some disagreements are method-internal
artifacts (one method reads a basis as canonical that another method
does not) rather than disagreements about the underlying model
computation.
